% TOP_PROBLEM: {"id": 3100001, "title": "How many colors is it necessary to use so that, if you paint every single point of the two-dimensional plane some color", "category_id": 2, "category": {"id": 2, "name": "combinatorics", "display_name": "Combinatorics", "description": "Counting problems, graph theory, discrete structures.", "slug": "combinatorics", "order_index": 2, "created_at": "2026-07-31T15:26:25.671Z"}, "status": "partially_solved"}
% FIRST_POSTED: null
% ENTRY_KIND: statement_only
% !TeX program = lualatex
\documentclass[11pt]{article}
\usepackage[margin=25mm]{geometry}
\usepackage{fontspec}
\setmainfont{Latin Modern Roman}
\usepackage{amsmath,amssymb,mathtools}
\usepackage{unicode-math}
\setmathfont{Latin Modern Math}
\usepackage{xurl,hyperref}
\hypersetup{colorlinks=true,urlcolor=blue}
\setcounter{secnumdepth}{0}
\setlength{\parindent}{0pt}
\setlength{\parskip}{5pt}
\setlength{\emergencystretch}{3em}
\newcommand{\R}{\mathbb R}
\newcommand{\N}{\mathbb N}
\newcommand{\Z}{\mathbb Z}
\newcommand{\Q}{\mathbb Q}
\newcommand{\C}{\mathbb C}
\newcommand{\D}{\mathbb D}
\newcommand{\F}{\mathbb F}
\newcommand{\sref}[2]{\hyperlink{source:#1:#2}{[S#2]}}
\newcommand{\cataloguescope}[1]{#1}
\begin{document}
% UnsolvedMath ID 3100001; source catalogue code AMR-030-0001
\setcounter{section}{3100000}
\section{How many colors is it necessary to use so that, if you paint every single point of the two-dimensional plane some color}\label{3100001}
{\small UnsolvedMath ID 3100001 · Combinatorics}
\subsection{Definitions and mathematical statement}

\paragraph{Shared notation.}$\mathbb N$, $\mathbb Z$, $\mathbb Q$, $\mathbb R$, and $\mathbb C$ denote the natural numbers, integers, rationals, reals, and complex numbers, respectively. All additional parameters and hypotheses are specified in the question.


\paragraph{Statement recorded in the supplied source.}
Determine the least number of colors for a coloring of all points of $\mathbb R^2$ in which points at distance one have different colors. For the closed disk $B_r(0)$, determine the corresponding chromatic number as a function of $r$. In particular, at what radius does it first reach four?
\par\smallskip\textit{Formulation references:} \sref{3100001}{1}, \sref{3100001}{2}.
\subsection{Short English statement}
Determine the least number of colors for a coloring of all points of $\mathbb R^2$ in which points at distance one have different colors. For the closed disk $B_r(0)$, determine the corresponding chromatic number as a function of $r$. In particular, at what radius does it first reach four?
\subsection{Sources}
\begin{itemize}
\item[\hypertarget{source:3100001:1}{S1}] ULAM.AI and the UnsolvedMath Contributors, UnsolvedMath: A Curated Collection of Open Mathematics Problems, record 3100001 (AMR-030-0001). Catalogue attribution under CC BY 4.0.
\url{https://huggingface.co/datasets/ulamai/UnsolvedMath}.
\item[\hypertarget{source:3100001:2}{S2}] Original problem formulation and mathematical context.
\url{https://people.math.sc.edu/cooper/combprob.html}.
\end{itemize}
\label{end:3100001}
\end{document}
